\chapter{\textbf{Domain} Analysis} \label{Analysis}
This section will cover an analysis of the case presented by the \acrshort{DEPA} in section 3.1 and what challenges it presents. From then it will delve into necessary knowledge required to develop a solution for the case, starting with some ground truths regarding potato agriculture in Denmark in section 3.2. Monitoring of potato crops will be performed using the Copernicus Satellite API, which is described in section 3.3 along with the satellite capabilities. Finally, the remote sensing data from the satellites can be used to calculate vegetation indices, which are explored in section 3.4 and the differences between the vegetation indices.

\section{Analysis of the stakeholder presentation} \label{Analysis of the stakeholder presentation}
During the case presentation with the representative from the \acrshort{DEPA}, it was clear that they wanted a more optimal way of checking inspection site candidates for potential pesticide overuse checks. Currently they employ a limited number of random agricultural soil samples on different sites each year\cite{MiljoeOgFoedevareministeriet}, which potentially results in pesticide overuse going unnoticed. To combat this, the case presented a solution to use satellite image data and using NDVI index over time, to see if there is an abnormal drop in NDVI index, which could be an indicator for pesticide overuse.
This presents some additional challenges, as different crop types are sprayed with pesticides for various reasons and at varying times during the year, depending on the type of crop. In accordance to this, the \acrshort{DEPA} had a specific use case regarding seed potatoes and the use of the pesticide Reglone\cite{reglone / case PP} that is used for pre-harvest crop desiccation. The proposed solution is to look for a drop in the NDVI index from the satellite imagery from seed potato fields, which would indicate that pesticide use for pre-harvest crop desiccation has been employed, and then to search for fields where the NDVI index drops off faster than the norm, which could be an indicator of overuse.

To be able to perform this task, the following section will delve into potato farming standards and practices and necessary knowledge.

\section{Potato knowledge} \label{Potato knowledge}
There are several potato types that being sown and harvested in Denmark. Most commonly, table/eating potatoes for consumption (da: spisekartofler), seed potatoes for planting (da: læggekartofler) and starch potatoes grown for potato starch production (da: melkartofler).
This is an umbrella term that encapsulates a variety of different potatoes in each potato type. Some varieties of potatoes have earlier harvests and/or could require different growth environments, in terms of soil nutrients and moisture. 

\subsection{Growth period} \label{Growth Period}
Potatoes are planted in late March and throughout April where the soil is around 8 degrees Celsius and the soil has the correct moisture levels\cite{LandbrugogFødevarer}. Before planting the field has been ploughed and is divided into trenches with approximately 75 cm between each trench, leaving space for machinery to drive by, without damaging the potatoes. In the trenches the potatoes are planted with proximately 20-33 cm distance between each other and covered with 6-8 cm layer of dirt\cite{LandbrugogFødevarer}. Once planted the potato takes anywhere from 2-6 weeks to germinate, depending of soil temperature, moisture and pH levels\cite{vegtotable_potato_germination}.

Harvesting season varies depending on the type and potato variety. Some table potatoes are harvested in June throughout July, while others are being harvested throughout August till October, which affects both size and taste of the table potatoes. Seed potatoes are harvested in the later months, August throughout October, to ensure they are large enough and carry enough nutrients for a healthy harvest the following year\cite{PPOplæg}.

\subsection{Pests and diseases}\label{Pests and diseas}
During the early growth period, before the potato seeds start sprouting out of the ground the field is bare. This enables weeds to get enough light to start growing, until the potato sprouts are tall enough to cover the ground, closing light off from the weeds which prevents their growth. Combating weeds in the early growth period can be done chemically using pesticides or mechanically using machinery\cite{LandbrugogFødevarer}.

Insect pests can attack the potatoes from both the sprouts and the roots, which negatively effects the yield and the quality of the potatoes. Whether pesticides are used to combat these pests is dependent on the severity of each case, usually determined by the amount of insects per potato plant\cite{InsectPests}.

Additionally, the potatoes can be inflicted by plant diseases, most notably potato blight. Potato blight is a fungus that, if goes untreated, causes the potato to rot and become inedible. To combat potato blight, the farmer can employ practices such as only planting healthy potatoes, hill the potatoes in relation to sprout growth to lessen the chance that the potato tuber gets infected and spray with an agent that prevents the fungus growth regularly\cite{PotatoBlight}. 

\subsection{Pre-harvest desiccation agents} \label{Pre-harvest desiccation agents}
There is a variety of pesticides. Herbicides, that focus on combating weeds or other plant matter, Insecticides, for killing insect pests and fungicides, to prevent fungal based plant sickness\cite{pesticiderUmbrellaterm}. And while a variety of different pesticides are used throughout the year to combat pests and diseases, the focus point presented by the case is the use of herbicides, specifically pre-harvest crop desiccation agents for seed potatoes in August throughout October.
The \acrshort{DEPA} stated that the monitoring of the pre-harvest desiccation agent Reglone was the focus point of the case\cite{oplægPP}. 

Reglone is a pre-harvest crop desiccation agent used on potatoes for storage (E.g.. seed potatoes). Reglone is a contact agent that is sprayed on the green sprouts growing out of the potato tuber and with the help of sunlight it creates hydrogen peroxide, wilting the plant. The agent is highly effective and dessication can be seen within 3-10 days. A maximum of two dosages is allowed and last dosage no longer than seven days prior to harvest\cite{Middeldatabasen}.

During the interview with the agricultural school, the two teachers interviewed mentioned that Reglone was hardly ever used due to regulations and most potato farmed had resorted to using milder pre-harvest desiccation agent, Mizuki\cite{interview}. Mizuki is also a contact agent and only used for potatoes for storage (E.g. seed potatoes). While the effect of Mizuki is not as immediately present as Reglone, it is highly effective and desiccates the plant after 14-21 days. Two dosages are allowed with a five day gap and the last application can be no longer than seven days prior to harvest\cite{Middeldatabasen}.

\subsection{Dangers of pre-harvest desiccation agents} \label{Dangers of pre-harvest desiccation agents}
Both Reglone and Mizuki are corrosive liquids and needs to be handled using gloves, masks and eye protection. Coming in contact with the agents causes skin irritation and needs to be washed off immediately. Breathing in the mist from the agents can cause irritation in the airways and in severe cases cause irreparable lung damage or even death\cite{Middeldatabasen}.
Besides being harmful to humans, regulations state that the agents must not be applied closer than 30 meters from aquatic environments (lakes, rivers, etc.) and 20 meters from § 3-areas, which includes wild vegetation\cite{Paragraf3}\cite{Middeldatabasen}. This is due to the agents being highly poisonous to aquatic wildlife and the stability of the native ecosystem in bodies of waters and has a damaging effect on wild vegetation\cite{Middeldatabasen}. It must also be applied further than 20 meters from roads and residential areas, in order to protect passer-by and residents\cite{Middeldatabasen}.

\subsection{Soil analysis after desiccant application}\label{Soil analysis after desiccant application}
Reglone is considered to be very persistent in soil and can be detected for years after application\cite{reglone_sds_2023}. Despite this, the bioavailability is minimal following the next years planting of seed potatoes, so the persistence of Reglone in the regulated amounts has a minimal or no effect on future harvests\cite{fao_diquat_2013}.

To verify whether a farmer has upheld the regulations in relation to pre-harvest desiccation agents, the \acrshort{DEPA} performs about 1000 random soil samples each year\cite{MiljoeOgFoedevareministeriet}. 

\section{Copernicus Satellite API}\label{Copernicus Satellite API}
The Copernicus Programme is operated by the European Union in collaboration with the European Space Agency (ESA)\cite{copernicus_about}. The Copernicus Programme aims to provide comprehensive, freely accessible data for the purpose of environmental management, disaster response, sustainable development and more. 

The Copernicus Programme has a series of satellites in orbit, but for the case presented by the \acrshort{DEPA} the focus will lie on the Sentinel-2 satellites, as their mission is focused on land based and agricultural monitoring. The Sentinel-2 satellites are equipped with a multispectral instrument (MSI) that covers 13 spectral bands. 
The bands are as follows\cite{sentiwiki_sentinel2}:
\begin{itemize}
    \item Band 1 : Coastal Aeresol (blue and violet)
    \item Band 2 : Blue
    \item Band 3 : Green
    \item Band 4 : Red
    \item Band 5, 6 \& 7 : Red Edge
    \item Band 8 : NIR (Near Infrared)
    \item Band 8A : Narrow NIR
    \item Band 9 : Water Vapour
    \item Band 10 : Cirrus (Cloud Coverage)
    \item Band 11 \& 12 : SWIR (Short-Waved Infrared)
\end{itemize}
The remote sensing data from the Sentinel satellites provides the necessary spectral ranges to calculate several vegetation indexes for satellite crop monitoring, to monitor the health of potato crops and the desiccation thereof\cite{softfarm_vegetation_indices}.

\section{Vegetation indices}\label{Vegetation indices}
A vegetation index is a calculation derived from combining specific spectral bands from remote sensing data, that reflects the condition of vegetation\cite{softfarm_vegetation_indices}. 
One of the conditions of plant health shows itself in the levels of chlorophyll in the plants. Chlorophyll is the green pigment in chloroplast cells in plants, that absorbs the sunlight necessary for photosynthesis. It also reflects green light, making leaves appear green\cite{byjus_chlorophyll}. Chlorophyll contains nitrogen, which in turn makes it possible to make a correlation between plant health and soil nutrient levels. If vegetation appears stressed during the growth period, it could be the result of low soil nutrients\cite{eos_ndre}.

There are several different vegetation indices and while the \acrshort{DEPA} suggested the use of NDVI for the use case, other vegetation indices may be worth exploring during testing.

\subsection{NDVI}\label{NDVI}
The Normalized Difference Vegation Index (NDVI) is one of the most commonly used vegetation indices, as it focuses on the density of vegetation. This allows for monitoring plant growth and plant health, as stressed vegetation is not as green and will reflect more visible red light and absorb more near infrared light, compared to healthy, greener vegetation\cite{softfarm_vegetation_indices}.

The formula for NDVI is as follows:
\begin{equation}
    \text{NDVI} = \frac{\text{NIR} - \text{Red}}{\text{NIR} + \text{Red}}
    \label{eq:ndvi}
\end{equation}

NDVI ranges from -1 to 1, where high values in the range of 0.6-1.0 would be dense healthy vegetation, ranges of 0.2-0.5 would be \sout{sparse} considered sparse vegetation, lower values of 0-0.2 would be considered non-vegetated areas and below 0 is either non-vegated surfaces, water, snow or clouds\cite{NDVIRanges}.

\subsection{EVI}\label{EVI}
The Enhanced Vegetation Index was created as an improvement of NDVI by factoring in areas with a high leaf area index (LAI). One of the issues that NDVI can present is that is more easily saturated (less sensitive where there is high biomass) and more sensitive to soil background in sparse vegetation. EVI corrects for this by using blue reflection to correct from background soil noise, in addition to reducing the effect of atmospheric effects, such as aerosol\cite{softfarm_vegetation_indices}.

The formula is as follows:
\begin{equation}
    \text{EVI} = 2.5 \times \frac{\text{NIR} - \text{Red}}{\text{NIR} + 6 \times \text{Red} - 7.5 \times \text{Blue} + 1} \label{eq:evi}
\end{equation}

EVI follows the same range as NDVI, but is more sensitive in areas with high biomass and may yield more sensitive readings\cite{EVIRanges}. Though it is more sensitive to bright and dark features, like clouds, buildings and water, which can result in anomalous EVI values\cite{softfarm_vegetation_indices}.

\subsection{GNDVI}\label{GNDVI}
The Green Normalized Difference Vegetation Index uses the same formula as NDVI but with the green spectrum instead of the red, to better monitor chlorophyll levels in the canopy of the plant, as chlorophyll pigment reflects green light. Because of this it is more suited towards spotting stressed or ageing vegetation\cite{softfarm_vegetation_indices}.

The formula is as follows:
\begin{equation}
    \text{NDVI} = \frac{\text{NIR} - \text{Green}}{\text{NIR} + \text{Green}}
    \label{eq:ndvi}
\end{equation}
GNDVI follows the same range as NDVI, but often have lower values in the same areas compared to NDVI (depending on crop type, canopy condition and growth stage), due to the green spectrum focusing on the chlorophyll concentration and not saturating like NDVI\cite{softfarm_vegetation_indices}.

\subsection{NDRE}\label{NDRE}
The Normalized Difference Red Edge Index is a method for monitoring chlorophyll levels in plants. Chlorophyll absorbs red light and reflects near-infrared light. In-between the red and near-infrared wavelengths is the red-edge, which describes the steep reflectance increase, as seen on figure (INDSÆT BILLEDE DU GEMTE I DISCORDEN)\cite{geopera_red_edge}.
Additionally, NDRE is designed for deep vegetation insight as red-edge wavelengths penetrate deeper into leaf layers and gets a more accurate picture of the overall plant health\cite{eos_ndre}.

The formula is as follows:
\begin{equation}
    \text{NDRE} = \frac{\text{NIR} - \text{RedEdge}}{\text{NIR} + \text{RedEdge}}
    \label{eq:ndre}
\end{equation}
NDRE follows the same range as NDVI, but generally has lower values, as it does not saturate on the canopy like NDVI.

\subsection{CVI}\label{CVI}
The Chlorophyll Vegetation Index, like the GNDVI, focuses on the concentration of chlorophyll in plants. It is most commonly used in the early and middle of crop cycles to determine plant health. If the plant lacks the proper nutrients chlorophyll-related stress can present itself, resulting in a lower CVI value\cite{softfarm_vegetation_indices}.

The formula is as follows:
\begin{equation}
    CVI = \left(\frac{NIR}{\text{Green}}\right)\left(\frac{\text{Red}}{\text{Green}}\right)
    \label{eq:cvi}
\end{equation}
CVI does not have a normalized values and can vary widely. Therefore plant health analysis through CVI is more context-dependent in terms of external factors (E.g. low soil nutrients). If low values are observed in certain regions, manual investigation is necessary to pinpoint the cause of the chlorophyll-related stress\cite{CVI}.
